\chapter{相关硬件基础}

\section{中央处理器}
\concept{CPU 的基本组成与工作流程}
\begin{points}
  \item CPU 的主要部件包括控制器、运算器、寄存器和高速缓存。
  \item CPU 的任务是执行指令，基本流程为：取指 $\rightarrow$ 译码 $\rightarrow$ 取操作数 $\rightarrow$ 执行指令。
  \item 影响 CPU 性能的因素包括主频、字长、指令集、高速缓存、核心数量、总线速度和处理技术。
\end{points}
\plain{硬件基础不用背太深，重点是知道哪些硬件机制支撑 OS：寄存器保存现场，PSW 保存状态，CPU 模式限制权限。}
\exam{常考 CPU 工作流程、单/多处理器、寄存器类别、PSW 中断位/状态位。}


\concept{单处理器与多处理器}
\begin{points}
  \item 单处理器系统只有一个运算处理器，同一时刻只能有一个进程在 CPU 上运行。
  \item 单 CPU 中，进程与进程只能并发不能并行；处理机与设备、设备与设备在硬件支持下可以并行工作。
  \item 多处理器系统有多个运算处理器，可分为对称多处理 SMP 和非对称多处理 AMP。
  \item SMP 中各处理器地位相同，共享内存和 I/O 系统；AMP 中主从处理器角色不同，主处理器负责调度与管理。
\end{points}
\plain{硬件基础不用背太深，重点是知道哪些硬件机制支撑 OS：寄存器保存现场，PSW 保存状态，CPU 模式限制权限。}
\exam{常考 CPU 工作流程、单/多处理器、寄存器类别、PSW 中断位/状态位。}


\concept{Flynn 分类}
\begin{points}
  \item SISD：单指令流单数据流，传统顺序单处理器。
  \item SIMD：单指令流多数据流，一条指令同时处理多个数据，如向量处理、GPU 中的某些模式。
  \item MISD：多指令流单数据流，实际很少见。
  \item MIMD：多指令流多数据流，多处理器/多核系统常见。
\end{points}
\plain{硬件基础不用背太深，重点是知道哪些硬件机制支撑 OS：寄存器保存现场，PSW 保存状态，CPU 模式限制权限。}
\exam{常考 CPU 工作流程、单/多处理器、寄存器类别、PSW 中断位/状态位。}


\section{寄存器、特权指令与处理器状态}
\concept{寄存器}
\begin{points}
  \item 寄存器是 CPU 内部高速存储单元，用于保存指令、地址、数据和状态信息。
  \item 常见寄存器包括程序计数器 PC、指令寄存器 IR、通用寄存器、栈指针 SP、程序状态字 PSW。
  \item 上下文切换时，操作系统需要保存当前进程的寄存器现场，并恢复下一个进程的现场。
\end{points}
\plain{硬件基础不用背太深，重点是知道哪些硬件机制支撑 OS：寄存器保存现场，PSW 保存状态，CPU 模式限制权限。}
\exam{常考 CPU 工作流程、单/多处理器、寄存器类别、PSW 中断位/状态位。}


\concept{特权指令与非特权指令}
\begin{points}
  \item 非特权指令可由用户程序直接执行，如普通算术、逻辑、数据移动指令。
  \item 特权指令只能由内核执行，如 I/O 指令、关中断、修改内存保护寄存器、修改页表、启动 DMA、设置定时器。
  \item 判断标准：凡是可能影响系统安全、影响其他进程、直接控制硬件资源的操作，通常属于特权操作。
\end{points}
\plain{普通程序不能直接碰全局危险资源；一旦需要 I/O、页表、中断等操作，就要进内核态。}
\exam{常考哪些指令是特权指令、系统调用是否引起模式切换、模式切换和进程切换区别。}


\concept{处理器状态：用户态与内核态}
\begin{points}
  \item 用户态下执行普通用户程序，只能执行非特权指令，不能直接访问受保护资源。
  \item 内核态下执行操作系统内核，可执行特权指令，访问核心数据结构。
  \item 从用户态到内核态的常见原因：系统调用、异常、外部中断。
  \item 从内核态返回用户态：中断/系统调用处理完成后，恢复用户进程现场并执行返回指令。
\end{points}
\plain{硬件基础不用背太深，重点是知道哪些硬件机制支撑 OS：寄存器保存现场，PSW 保存状态，CPU 模式限制权限。}
\exam{常考 CPU 工作流程、单/多处理器、寄存器类别、PSW 中断位/状态位。}


\concept{PSW 程序状态字}
\begin{points}
  \item PSW 保存 CPU 的运行状态信息，包括条件码、中断屏蔽位、处理器模式位等。
  \item 模式位用于区分用户态和内核态，是硬件支持操作系统保护机制的重要基础。
  \item 中断发生时，硬件通常会保存 PC/PSW 等关键信息，以便中断处理完成后回到原程序继续执行。
\end{points}
\plain{硬件基础不用背太深，重点是知道哪些硬件机制支撑 OS：寄存器保存现场，PSW 保存状态，CPU 模式限制权限。}
\exam{常考 CPU 工作流程、单/多处理器、寄存器类别、PSW 中断位/状态位。}


\section{中断处理}
\concept{中断的基本概念}
\begin{points}
  \item 中断是 CPU 暂停当前程序执行，转而执行相应处理程序，处理完成后再返回原程序的机制。
  \item 中断使 CPU 不必忙等 I/O 完成，是实现多道程序设计和 CPU/I/O 并行工作的关键。
  \item 中断来源包括外部设备中断、定时器中断、异常和系统调用陷入。
\end{points}
\plain{中断就是硬件或软件打断 CPU 当前执行，让 OS 先处理更紧急或必须处理的事件。}
\exam{常考中断处理步骤、中断与异常/系统调用区别、为什么中断能提高 CPU 与 I/O 并行度。}


\concept{中断处理过程}
\begin{points}
  \item 设备或内部事件提出中断请求。
  \item CPU 在适当时刻响应中断，保存断点和现场，切换到内核态。
  \item 根据中断向量找到对应的中断服务程序。
  \item 中断服务程序处理中断事件，例如读取设备状态、唤醒等待进程。
  \item 恢复被中断程序现场，返回原程序或转入调度程序。
\end{points}
\plain{中断就是硬件或软件打断 CPU 当前执行，让 OS 先处理更紧急或必须处理的事件。}
\exam{常考中断处理步骤、中断与异常/系统调用区别、为什么中断能提高 CPU 与 I/O 并行度。}


\concept{中断与异常、系统调用}
\begin{points}
  \item 外部中断通常来自 CPU 外部设备，如磁盘、网卡、键盘。
  \item 异常来自当前指令执行过程，如除零、非法指令、缺页异常。
  \item 系统调用是用户程序有意执行陷入指令，请求内核服务。
  \item 三者都会导致从用户态进入内核态，但触发原因不同。
\end{points}
\plain{人用命令/GUI 找 OS，程序用系统调用找 OS；危险操作必须让内核代办。}
\exam{常考命令接口、程序接口、GUI 的分类，以及系统调用是否属于程序接口。}


\section{操作系统的硬件支持}
\concept{内存保护}
\begin{points}
  \item 为防止进程访问非法内存，硬件需要支持地址变换和地址合法性检查。
  \item 简单系统中可用基址寄存器和界限寄存器检查地址是否越界。
  \item 分页/分段系统中由 MMU 根据页表或段表完成地址映射和保护检查。
\end{points}
\plain{把这个知识点放回本章主线理解：它回答的是 OS 为了管理资源、隐藏硬件、保证并发正确性而采用的某个对象、规则或步骤。}
\exam{常考选择/填空/简答，让你判断概念归属、比较相近概念，或按“定义-作用-机制-易错点”展开。}


\concept{定时器}
\begin{points}
  \item 定时器用于防止用户程序长期独占 CPU。
  \item 操作系统设置时间片，到时后定时器中断发生，CPU 转入内核执行调度程序。
  \item 分时系统和抢占式调度离不开定时器支持。
\end{points}
\plain{把这个知识点放回本章主线理解：它回答的是 OS 为了管理资源、隐藏硬件、保证并发正确性而采用的某个对象、规则或步骤。}
\exam{常考选择/填空/简答，让你判断概念归属、比较相近概念，或按“定义-作用-机制-易错点”展开。}


\concept{I/O 保护}
\begin{points}
  \item 用户程序不能直接执行 I/O 指令，必须通过系统调用请求内核完成。
  \item 这样可以统一设备分配、访问权限、缓冲、错误处理，防止多个进程直接抢设备。
  \item 因此 I/O 指令属于典型特权指令。
\end{points}
\plain{I/O 部分就是解决 CPU 快、设备慢、设备复杂且共享的问题。DMA/通道减少 CPU 搬运，缓冲平滑速度差，SPOOLing 把独占设备虚拟成共享设备。}
\exam{常考设备分类、四种 I/O 控制方式、DMA 与中断/通道区别、缓冲和 SPOOLing 的作用。}


\section{计算机启动与虚拟机}
\concept{计算机启动过程}
\begin{points}
  \item 上电后 CPU 从固定地址开始执行固件中的引导程序。
  \item 固件进行硬件自检和基本初始化，找到可启动设备。
  \item 引导程序把操作系统内核装入内存。
  \item 内核初始化数据结构、设备驱动、内存管理、进程管理，创建第一个用户态进程或系统服务。
\end{points}
\plain{开机不是直接进 OS，而是固件先跑，找到引导程序，再把 OS 内核装入内存。}
\exam{常考 BIOS/UEFI、引导程序、内核加载的大致顺序。}


\concept{虚拟机}
\begin{points}
  \item 虚拟机把底层物理机器抽象成多台逻辑机器，使多个 OS 能在同一物理平台上运行。
  \item 虚拟机监控器负责 CPU、内存、设备等资源的虚拟化与隔离。
  \item 直觉：普通 OS 把裸机虚拟成“适合应用使用的机器”；虚拟机监控器把裸机虚拟成“多台可运行 OS 的机器”。
\end{points}
\plain{不是资源真的变多，而是 OS 用映射、换入换出、排队等办法，让用户感觉自己有一份独立资源。}
\exam{常考虚拟处理机、虚拟存储器、虚拟设备的例子，以及虚拟性和资源复用的关系。}


